\lecture{9}
\title{Maximum Flow (Part I)}

\begin{document}

\begin{problem}{Max Flow}
    Consider a directed graph $G = (V, E)$ with a source vertex $s$ and a sink vertex $t$.  Every (directed) edge has a capacity $c: E \to \mathbb{R}^{\geq 0}$, extended to all of $V \times V$ by $c(u, v) = 0$ when $(u, v) \not \in E$.
    
    We'd like to assign a flow to every edge such that:
    \begin{enumerate}
        \item For any $e \in E$, the flow across $e$ is at most the capacity of $e$.
        \item There is zero net flow into or out of any vertex other than $s$ or $t$.
    \end{enumerate}
    Formally, let $f: V \times V \to \mathbb{R}$ such that $f(u, v)$ is the (signed) flow from $u$ to $v$, with $f(u, v) = -f(v, u)$ and $f(u, v) \leq c(u, v)$ for all $u, v \in V$ (so $f(u, v) = 0$ if neither $(u, v)$ nor $(v, u)$ is in $E$).

    The net flow $|f|$ is defined by $\sum_{v \in V} f(s, v)$.  Find the maximum possible total flow.
\end{problem}

\begin{center}
    \frame{\includegraphics[width=0.45\textwidth]{lecture-09/flow-example}}
\end{center}

\begin{claim} % Flow Decomposition Lemma
    Any flow assignment can be decomposed as the linear combination of \emph{paths from $s$ to $t$} and \emph{cycles}.
\end{claim}

\begin{proof}
    Just induct.  You never get ``stuck'' because of the ``zero net flow'' constraint, which is invariant. ~ $\blacksquare$
\end{proof}

There's an easy upper-bound on the max flow: the minimum cut.

\textbf{Definition (Min Cut).} Cut the directed graph via $V = S \sqcup T$, with $s \in S$ and $t \in T$.  We define:
\begin{itemize}
\item The total cut $c(S)$ is the sum of the (nonnegative) capacities $c(a, b)$ for all $a \in S$ and $b \in T$.
\item The total flow $f(S)$ is the sum of the (not necessarily nonnegative) flows $f(a, b)$ for all $a \in S$ and $b \in T$.
\end{itemize}
Evidently, $f(S) \leq c(S)$.

\begin{claim} % Weak Duality
    Consider any flow assignment $f$.  Given any two cuts $S_1 \sqcup T_1$ and $S_2 \sqcup T_2$, we have $f(S_1) = f(S_2)$.
    
    (By corollary, for any cut $S \sqcup T$, we have $f(S)$ is the net flow of $f$.)
\end{claim}

\begin{proof}
    Because the divergence is zero.  (Formally, use the decomposition of flow into paths and cycles.) ~ $\blacksquare$
\end{proof}

Be careful, though---the above does \emph{not} say that all cuts have the same cut.  It only says that all cuts have the same \emph{flow}---that is, the same measure with respect to edge weights with divergence $0$ everywhere.

\begin{boxed-text}
    \textbf{Claim.} The maximum possible flow of a graph equals the minimum possible cut.
\end{boxed-text}

Here's the key idea.

\textbf{Definition (Residual Network).} Given an assigned flow $f$ to a graph $G = (V, E)$, the \textbf{residual network} consists of $G_f = (V, E_f)$ with edge weights $c_f(u, v) := c(u, v) - f(u, v)$, and edges $E_f := \{(u, v) \in V \times V \mid c_f(u, v) > 0\}$.

Importantly, some edges in $G_f$ will have \emph{greater} weight than the corresponding edge in $G$, because flow along an edge $u \to v$ increases $c_f(v, u)$.

\begin{center}
    \frame{\includegraphics[width=0.75\textwidth]{lecture-09/residual-network}}
\end{center}

Think of this as reflecting our ability to ``undo our mistakes''.

\textbf{Definition (Augmenting Path).} An \textbf{augmenting path} is a path $s \to \dots \to t$ in the residual network $G_f$.

\begin{claim}
    The max flow is zero if and only if there is a cut of zero.
\end{claim}

\begin{proof}
    If there's a zero cut, there's obviously zero flow.  If there is zero possible flow, that means there is no path from $s$ to $t$ of positive capacity, so the cut defined by $S := \{v \in V \mid \text{there is a positive-capacity path from } s \text{ to } v\}$ works. ~ $\blacksquare$
\end{proof}

It turns out that the zero case is all we need!  (Like proving JNF\dots)

\begin{claim}
    The maximum possible flow is equal to the minimum possible cut.
\end{claim}

\begin{proof}
    Consider an optimal flow assignment $f$.  Then $G_f$ must have zero maximum possible flow, or else $f$ would not be optimal.  This means $G_f$ has a cut $S \sqcup T$ of zero.

    \begin{quote}
        \textbf{Claim.} For any graph $G$, any cut $S$, and any flow $f$, we have $c_f(S) = c(S) - f(S)$.

        \emph{Proof:} Pretty obvious.  Do the algebra if you want. ~ $\square$
    \end{quote}    
    When there is a cut $S \sqcup T$ of value zero in $G_f$, we have $f(S) = c(S)$.  Lower bound meets upper bound. ~ $\blacksquare$
\end{proof}

More generally, the following is true.

\begin{claim}
    Given any network $G$ and any flow $f$, the following statements are equivalent.
    \begin{enumerate}
        \item We have that $f$ is a max flow.
        \item We have that $|f| = c(S)$ for some cut $S$.
        \item We have that $G_f$ has no augmenting paths.
    \end{enumerate}
\end{claim}

\begin{proof}
    Omitted to be DRY.  (I am not writing a helper proof.) ~ $\blacksquare$
\end{proof}

And so the most naive algorithm actually just works.

\begin{solution}{Ford-Fulkerson}
    Repeatedly find augmenting paths $P := s \to \dots \to t$, then recompute the residual network each time.  Repeat until no more augmenting paths exist.  This is optimal.
\end{solution}

\end{document}